% reconfiguration.tex -- Methodology 3.4: reconfiguration (rewritten 12 Oct 2026 to the thesis plan of 11 Oct:
% notice, re-plan, redistribute, fall back; the three-equation statics; the support margin with its qualifier; the
% constants table; the admission contract; the state machine). 17 Oct (v3): the join is noticed by the co-motion
% confirmation (the default since 14 Oct, proof_rules comotion; card docs/experiments/2026-10-15_attach_comotion_rig.md
% rev 4, Tejen's review); the old contact-and-lift proof is history. Plates clockwise (angle = 30 x plate).
% Code: the planner's detach mode (detection, resize), drone_magnet (contact proof, attachment manager, radio mux),
% the take-back in the tracker; the collaborator's stack carries the lift from contact.

\subsection{Reconfiguration}
\label{sec:reconfig}

A reconfiguration changes the number of UAVs carrying the ring in the air. Both directions are handled by the same
four steps: the planner \emph{notices} the change from an observation, \emph{re-plans} by switching to a solver
built for the new team, \emph{redistributes} the ring's weight across it, and \emph{falls back} to a simple exit
when a step fails. Figure~\ref{fig:reconfig_states} shows the two lanes and their exits; the constants every step
uses are in Table~\ref{tab:constants}, each with the reason for its value.

\begin{figure}[htbp]
\centering
\begin{tikzpicture}[>=Stealth, font=\small,
    st/.style={draw, thick, rounded corners=3pt, align=center, minimum height=0.8cm, minimum width=2.2cm, font=\small},
    ours/.style={st, fill=blue!6, draw=blue!50!black},
    his/.style={st, fill=orange!10, draw=orange!60!black},
    exit/.style={st, fill=gray!14, draw=gray!60},
    ev/.style={font=\scriptsize, align=center}]
  % leave lane (top)
  \node[ours] (carry)   at (0,0)      {Carry, $N$ UAVs\\ \scriptsize planner owns the ring};
  \node[ours] (resize)  at (4.8,0)    {Resize to $N-1$\\ \scriptsize warm start, free tensions};
  \node[ours] (carry3)  at (9.6,0)   {Carry, $N-1$\\ \scriptsize statics sets the shares};
  \node[ours] (freed)   at (4.8,-1.7) {Freed UAV\\ \scriptsize steps out 0.5\,m};
  \node[exit] (dland)   at (9.6,-1.7){Lands by itself\\ \scriptsize disarmed alone at rest};
  \draw[->, thick] (carry) -- node[ev, above] {DETACH, or\\ rod stretch} (resize);
  \draw[->, thick] (resize) -- (carry3);
  \draw[->, thick] (resize) -- node[ev, right] {the leaver} (freed);
  \draw[->, thick] (freed) -- (dland);
  % join lane (middle)
  \node[his]  (mission) at (0,-3.6)   {Newcomer, free\\ \scriptsize collaborator's controller};
  \node[his]  (contact) at (4.8,-3.6) {Descent to contact\\ \scriptsize then a small lift};
  \node[ours] (proof)   at (9.6,-3.6){Confirmed by mocap\\ \scriptsize tip moves with the plate};
  \node[ours] (takeback) at (4.8,-5.3){Take-back\\ \scriptsize moved clear, lands alone};
  \node[ours] (fold)    at (9.6,-5.3){Mux switch, fold in\\ \scriptsize seed 0.1\,N, resize to $N+1$};
  \draw[->, thick] (mission) -- node[ev, above] {transit\\ to the plate} (contact);
  \draw[->, thick] (contact) -- node[ev, above] {contact} (proof);
  \draw[->, thick] (proof) -- node[ev, right] {confirmed} (fold);
  \draw[->, thick, dashed] (contact) -- node[ev, right] {give-up} (takeback);
  \draw[->, thick, dashed] (proof.south west) -- node[ev, sloped, above, pos=0.45] {not confirmed} (takeback.north east);
  \draw[->, thick] (fold.east) -- ++(0.55,0) |- node[ev, pos=0.25, right] {carry, $N+1$} (carry3.east);
  % exits (bottom)
  \node[exit] (drop) at (0,-7.0)    {Drop and land\\ \scriptsize second loss, or tilt $>$ 30$^\circ$};
  \node[exit] (cut)  at (4.8,-7.0)  {Tilt cut\\ \scriptsize newcomer's motors off};
  \node[exit] (land) at (9.6,-7.0) {LAND\\ \scriptsize every UAV, from any phase};
  \draw[->, thick, dashed] (takeback) -- node[ev, right] {stuck, tilting} (cut);
\end{tikzpicture}
\caption[Leave and join lanes and their exits]{The leave lane (top), the join lane (middle) and the exits
(bottom). Blue: this work's planner and trackers; orange: the collaborator's controller; grey: an exit that ends
with every UAV owned and on the floor. A leave is noticed from the DETACH command or from the leaver's rod reading
longer than itself; a join from a contact and, after the collaborator's lift, the magnet tip seen moving with the
plate on a taut tether, from motion capture alone. Every exit is reachable from every state: LAND and the drop rule
are not drawn into each.}
\label{fig:reconfig_states}
\end{figure}

\subsubsection{Notice}

\paragraph{A leave.} An \emph{announced} leave is the DETACH command: the leaver's magnet is switched off on its
radio channel and the planner resizes on the same tick, or, for a \emph{ramped} release, its planned tension is run
down over a short ramp before the magnet is released, so that the survivors take up the share before the rod goes.
An \emph{unannounced} leave is a magnet switched off without the planner being told, which is how a release the
planner does not expect is made on the rig and in the twin. The planner notices it from the rod: a rigid rod on its
plate cannot measure longer than itself, so when the leaver's rod, measured from motion capture between its pivot
and its plate point, reads more than 0.06\,m over its length for two consecutive ticks, the planner drops that UAV
from its problem. Nothing else is needed: no force sensor, no message from the leaver.

\paragraph{A join.} A join is noticed in two steps, both read from motion capture and nothing else, with the
magnet tip expressed in the plate's own frame. A \emph{contact} is the tip inside the catch radius of the plate
point (24\,mm on the rig), between 15\,mm below and 12\,mm above the plate's face, and slower than 0.03\,m/s for
0.3\,s. A contact is only a candidate: a tip resting on the plate and a magnet that has caught look alike. The
\emph{confirmation} uses the collaborator's small lift: 1.0\,s after its controller signals the lift done, the catch
is confirmed if the tip has stayed within 6\,mm of its contact point while the plate has moved with it (a range of
5\,mm or a rise of 3\,mm), the tip is within 10\,mm of the plate's face, and the tether reads taut against the free hang measured on the
same flight (within 3\,mm). A separation of 45\,mm for 0.12\,s fails it, and no hold spans missing data: a gap of more than 0.1\,s
restarts the contact and the confirmation. In the twin, whose stiff ring hardly moves and whose magnet catches at the
instant of contact, the plate-motion condition is waived and the face window is 15\,mm instead of 10. Until 14 October the proof
asked instead for the tip-to-plate vector to stay within 3\,mm while the plate rose 3\,mm; on the rig it rejected the
one catch of 14 October on its taut check, and this co-motion confirmation replaced it.

\subsubsection{Re-plan}

The planner keeps flying the ring and re-solves for the new team, which two properties of the formulation make
possible. The attachment points $\boldsymbol{\rho}_i$ and rod lengths $l_i$ are runtime parameters of the model, not
constants compiled into the solver, so only the team size fixes the problem's dimensions, and one solver per team
size, built in advance and primed twice on the measured state before it flies, covers any layout. And the
survivors keep their own attachment points, plates 3, 7 and 11 when the carry layout loses plate 9: the planner's
nominal tensions come from a moment balance at the actual plates, so the ring's reference stays level, whereas a
re-spaced even ring would make the moment balance wrong by construction. At the event the planner swaps to the
solver for the new size, sets the attachment points and rod lengths (for a newcomer, the plate it caught, recorded
in the ring frame at the confirmation), and warm-starts the new problem from the measured state with the tensions
and rod directions free at its first node (Section~\ref{sec:controller}); the first solve redistributes the load,
and there is no separate hand-off phase. An announced leave is refused during LAND and when it would leave fewer
than three UAVs, since a rigid ring on two rim points is a pendulum about the line between them; an unannounced
loss that leaves fewer than three is a fault, which drops the ring (Section~\ref{sec:safety}).

\subsubsection{Redistribute}

\paragraph{What statics fixes.} With three carriers the vertical shares are not the planner's to choose. Take the
ring's centre of mass as the origin of its frame and the survivors' plates at $\boldsymbol{\rho}_i = r(\cos\alpha_i,
\sin\alpha_i, 0)$; with $f_i$ the vertical component of rod $i$'s pull on a level ring, the force balance and the
two moment balances about the horizontal axes are three equations in three unknowns,
\begin{equation}
    \sum_{i=1}^{3} f_i = m_L g, \qquad \sum_{i=1}^{3} f_i \cos\alpha_i = 0, \qquad \sum_{i=1}^{3} f_i \sin\alpha_i = 0,
    \label{eq:statics}
\end{equation}
whose solution is unique whenever the three plates do not lie on a line: each share is the weight times the
barycentric coordinate of the centre of mass in the survivors' triangle. The rods' horizontal components, which the
planner does set through the rod directions, balance one another and make no moment about the centre of mass when
the centre lies in the plates' plane \cite{michael2011cooperative, tognon2018internal}; on the
ring the plates sit 2\,cm above the centre of mass, so the shares of \eqref{eq:statics} hold to within a few per
cent. After a leave the planner's first solve therefore hands the leaver's share to the survivors as
\eqref{eq:statics} says (Figure~\ref{fig:detach}). A level ring is possible only when every $f_i$ is positive, which
is the condition of the next paragraph.

\paragraph{Support margin.} A rigid payload can hang level from three or more rim points only if its centre of mass
lies inside the polygon of the attachment points \cite{michael2011cooperative}. For plates on a ring of radius $r$
the nearest edge of the survivors' polygon is the chord across their largest angular gap $g_{\max}$, and the
distance from the centre to that chord,
\begin{equation}
    m = r \cos\!\big(g_{\max}/2\big),
    \label{eq:margin}
\end{equation}
is the support margin: positive inside, zero on the edge, negative outside. On the 0.225\,m plate radius a gap of
120, 135, 150, 165 or 180$^\circ$ leaves a margin of 112, 86, 58, 29 or 0\,mm. Removing one carrier from an even
ring of $N$ leaves a largest gap of $4\pi/N$, which keeps the centre inside only for $N \ge 5$; an even ring of
four sits exactly on the boundary after any loss, which is why the four-UAV flights use plates 3, 7, 9 and 11,
where the loss of the carrier on plate 9 leaves an even triangle with a 112\,mm margin. The margin is the lever the
results chapter varies in the twin (Section~\ref{sec:res_leave}).

\paragraph{Which carrier may leave, and what control still sets.} The condition also says which carriers a layout
can afford to lose, and it can be checked before a flight from the plate angles alone. Removing one carrier merges
the two angular gaps on either side of it, and the four gaps of a four-carrier layout sum to $360^\circ$, so some
pair of adjacent gaps sums to at least $180^\circ$: no four-carrier layout on the ring can lose an arbitrary
carrier and still hang level, and a layout admits at most two leavers. The carry layout admits exactly one, the
carrier on plate 9; losing the carrier on plate 7 or plate 11 puts the centre on an edge, and losing the one on plate
3 puts it 112\,mm outside. The layout on plates 3, 6, 9 and 11 admits the same leaver with a $150^\circ$ gap and a
58\,mm margin. A team that must tolerate the loss of any one carrier therefore needs five, and a four-carrier layout
is designed for a known leaver, which is what the rejoin mission does.
% The numbers: tools/admissibility_map.py (results/analysis/admissibility/), from metrics.cog_margin and the
% planner's balanced_tensions; 0 of the 495 four-plate layouts admits every leaver (11 Oct).
Once the leaver has gone, what control can still set is bounded by the same statics. With three survivors no
controller can change how the ring's weight is shared. What the planner does set is the margin, through the choice
of layout and leaver before the flight; when and how fast the survivors take up the step, since an announced leave
can be ramped while an unannounced one is taken up at the detection; the rods' directions, which fix the horizontal
balance and where the survivors fly; and the ring's position and height through the transient. With four or more
carriers the tensions have a one-parameter family of solutions and the planner takes the split closest to equal
shares, so there a level ring can be held with more than one distribution.

\paragraph{The ramp and the seed.} Where the planner is told of a change it can shape the step. For a ramped
release the leaver's planned tension is run down and the survivors' run up over a short ramp before the magnet is
released, in place of a step at the resize tick; the first version of the ramp, which ran the leaver's tension
down but let the survivors' step, is the negative that defined it. For a join the newcomer is seeded at 0.1\,N and
every tension is blended to the balanced solution for the new geometry, as a dock blends thrust across its
controller switch \cite{sugihara2023trady}: on the earlier plant a newcomer started at a full share pulled the ring
with a lever about the incumbents' rod apex and tilted it at the weld (Appendix~\ref{sec:app_proximity_weld}).

\subsubsection{Fall back}

\paragraph{The freed UAV.} After a release the freed UAV steps 0.5\,m out radially at its height, clear of the ring
and the survivors, and lands by itself while the survivors carry on: its height reference is pressed below its own
resting height until it is there, and it is then disarmed alone, without waiting for the fleet
(Section~\ref{sec:res_freed}). Touchdown for the fleet is judged on the survivors only, so a freed UAV still
descending elsewhere does not hold their references down, and the fleet is declared landed only when it is down
too; the rod feedforward is switched off once the ring rests, because its sideways component pushes a grounded
survivor across the floor.

\paragraph{The take-back.} A newcomer whose contact is not confirmed within its timeout, or whose controller gives
up, is taken back: this work's tracker latches onto its measured position and velocity, ramps its reference at
0.25\,m/s$^2$ so that the latch makes no step, lifts it clear of the rim and the carriers, moves it out and lands it
by itself. On the circle the hold moves with the ring in the ring's planned frame, so that an orbiting formation
cannot sweep a carrier into the newcomer's rod; where the newcomer is low and far out, it lands where it stops. On
the rig the take-back first rises only 5\,cm above the tether's slack and waits for the tip to leave the plate, so
that a magnet that has caught after all is not torn off; if the plate rises instead, the magnet is switched off
again and the fleet lands. The tilt cut, the drop rule and LAND are the exits of Section~\ref{sec:safety}.

\subsubsection{The admission contract}
\label{sec:admission}

The join is a contract between two controllers, and this work's side of it is written so that nothing on the
newcomer's side has to be trusted. The collaborator's controller, which flies the newcomer to the ring, must:
descend onto its plate to contact with the magnet on (on the rig it waits at the bottom with the magnet on until a
contact is seen); lift by a small set amount and hold; signal that the lift is done; and give up its authority when
the catch is confirmed, or give up and go home. This work must: observe the contact and the lift from motion capture
only (the tip in the plate's frame, its drift from the contact point, the plate moving with it and the taut
tether); switch the newcomer's radio to its own tracker on the same tick as the confirmation; fold the newcomer in
from 0.1\,N; and take it back if nothing is confirmed. The multiplexer in front of the newcomer's radio forwards the
collaborator's commands until the confirmation, switches only onto a command that is armed and above idle throttle
(a switch onto an idle stream cut the newcomer's motors at every weld in early tests), and latches; an abort latches
it to a disarmed idle command until it is restarted. The fold-in records the plate in the ring frame, switches the
planner to the solver for $N+1$, and the trajectory continues. Section~\ref{sec:res_join} scores the contract on what
each side costs the ring: the time from contact to confirmation, the ring's tilt rise and dip through the admission,
the step in the newcomer's commands at the switch, and the clearances.

\paragraph{The rejoin mission.} The full mission is assembled with the collaborator's mission planner. Four UAVs
carry the ring in hover or round the circle; the one on plate 9 is released and flies the collaborator's errand
(picking up an object with its magnet and dropping it into the ring) and comes back to its plate; it is admitted
under the contract above, folded in, and the fleet lands. On the rig the collaborator's controller does every
attach down to contact, since the magnet catches only on the plate; this work's own approach, which steered the
measured magnet tip and welded on proximity, flies in the twin only (Appendix~\ref{sec:app_proximity_weld}).

\begin{table}[htbp]
\centering
\small
\caption[Reconfiguration constants and their reasons]{The constants of the reconfiguration and the reason for each
value: a noise floor times a factor, or a physical bound. [Bracketed reasons wait for the offline sweeps over the rig
recordings.]}
\label{tab:constants}
\begin{tabularx}{\textwidth}{>{\raggedright\arraybackslash}p{3.6cm}>{\raggedright\arraybackslash}p{2.4cm}>{\raggedright\arraybackslash}X}
\toprule
Constant & Value & Reason \\
\midrule
Rod stretch for an unannounced leave & 0.06\,m, two ticks & A rod on its plate cannot read longer than itself; the
threshold sits above the scatter of a rod length measured from motion capture, and two ticks reject a single bad
pose. [False alarms over the rig recordings and delay against threshold from the rig releases.] \\
Contact & catch radius 24\,mm (rig); face $-15$ to $+12$\,mm; under 0.03\,m/s for 0.3\,s & The magnet catches only on
the plate: the tip must be over the plate's face and at it, and still for 36 motion-capture samples, which excludes a
pass over the plate. \\
Confirmation, 1.0\,s after the lift & tip within 6\,mm of its contact point; plate range 5\,mm or rise 3\,mm; taut
within 3\,mm & A caught tip moves with the plate and a missed one is left behind by it; the free hang is measured on
each flight, so the taut check reads this flight's tether. \\
Confirmation fails & 45\,mm for 0.12\,s; any data gap over 0.1\,s restarts & A tip that leaves the plate has not
caught; no hold is judged across missing poses. \\
Newcomer's seed tension & 0.1\,N & Small against the newcomer's balanced share (1.6\,N on plate 9 of the carry
layout), so the fold-in's blend and not the weld sets the step. \\
Step-out after a release & 0.5\,m radially & Clears the rim (0.255\,m) and the survivors' rods from the leaver's
plate. \\
Take-back ramp & 0.25\,m/s$^2$ & A quasi-static tilt of 1.5$^\circ$ on the newcomer, under the 2$^\circ$ bar on its
attitude step, and it stops 0.34\,m/s in 1.4\,s. \\
Take-back clearances & 0.15\,m over the carriers; 0.3\,m under the ring; 0.9\,m from the ring's centre & The
newcomer's body and rod stay a rotor disc clear of the carriers; a low newcomer lands where it is rather than
rising through the formation. \\
Ring tilt that drops the ring & 30$^\circ$ & Beyond the planner's near-level model and about where a rod's pull turns
across its plate; where the controller gives up, not a value to design to (Section~\ref{sec:safety}). \\
Fewest carriers & 3 & Two rim points make the ring a pendulum about the line between them. \\
\bottomrule
\end{tabularx}
\end{table}
